\cleardoublepage

% concept04.tex
\section{Factors and Multiples}

Some numbers fit together cleanly. Others do not. Factors and multiples explain
why. A factor is a number that divides evenly. A multiple is what you get when
you multiply. These two ideas help us break numbers apart and build them back up.

If you list the factors of \(12\), you get \(1, 2, 3, 4, 6,\) and \(12\). Each
divides \(12\) without a remainder. The multiples of \(5\) are \(5, 10, 15, 20,\)
and so on. These ideas may seem mechanical, but they are powerful. They lay the
groundwork for division, fractions, and algebra.

\blockquote{A factor\index{Factor} is a number that divides evenly. 
A multiple\index{Multiple} is what you get when you multiply.}

\subsection{What are Factors and Multiples?}

As you continue your journey in mathematics, you'll often need to break numbers
down or build them up. This is where factors and multiples come in. A
\textbf{factor} is a number that divides another number exactly, leaving no
remainder. For example, the factors of \(12\) are \(1, 2, 3, 4, 6,\) and \(12\),
because each of these divides \(12\) evenly. On the other hand, a
\textbf{multiple} is what you get when you multiply a number by an integer. For
instance, the multiples of \(5\) are \(5, 10, 15, 20,\) and so on.

Let's see how this works with a few examples. If you want to find all the
factors of \(24\), you list out \(1, 2, 3, 4, 6, 8, 12,\) and \(24\). If you're
asked for the first \(6\) multiples of \(8\), you simply multiply: \(8, 16, 24,
32, 40,\) and \(48\). Sometimes, you might wonder if one number is a factor of
another. For example, is \(9\) a factor of \(81\)? Since \(81 \div 9 = 9\) with
no remainder, the answer is yes. Understanding factors and multiples is crucial
for topics like division, fractions, and finding patterns in numbers.



% \cleardoublepage
\subsection{Short Question (English)}
List all the factors of 18.

\fullpageimage{images/q1_4_1-eng.jpg}

\newpage
\subsection{Short Question (Korean)}
List all the factors of 18.

\fullpageimage{images/q1_4_1-kor.jpg}

\newpage
\subsection{Short Question (Answer)}
List all the factors of 18.

\textbf{Answer:} \textbf{1, 2, 3, 6, 9, 18}

% \cleardoublepage
\subsection*{Long Question (English)}
What are the first six multiples of 7? Explain how you found them.

\fullpageimage{images/q1_4_2-eng.jpg}

\newpage
\subsection*{Long Question (Korean)}
What are the first six multiples of 7? Explain how you found them.

\fullpageimage{images/q1_4_2-kor.jpg}

\newpage
\subsection*{Long Question (Answer)}
What are the first six multiples of 7? Explain how you found them.

\textbf{Answer:}\\
7, 14, 21, 28, 35, 42 — Multiply 7 by 1 through 6.

% \cleardoublepage
\subsection{Challenge Question (English)}
Is 9 a factor of 81? Is 81 a multiple of 9? Explain.

\fullpageimage{images/q1_4_3-eng.jpg}

\newpage
\subsection{Challenge Question (Korean)}
Is 9 a factor of 81? Is 81 a multiple of 9? Explain.

\fullpageimage{images/q1_4_3-kor.jpg}

\newpage
\subsection{Challenge Question (Answer)}
Is 9 a factor of 81? Is 81 a multiple of 9? Explain.

\textbf{Answer:}\\
Yes. $81 \div 9 = 9$ — so 9 is a factor, and 81 is a multiple.
